\documentclass[11pt]{article}
\usepackage[margin=0.88in]{geometry}
\usepackage{amsmath,amssymb,amsthm,mathtools}
\usepackage{graphicx}
\usepackage[hidelinks]{hyperref}
\usepackage{xcolor}
\usepackage{microtype}

\newtheorem{theorem}{Theorem}
\newtheorem{lemma}[theorem]{Lemma}
\newtheorem{proposition}[theorem]{Proposition}
\newtheorem{corollary}[theorem]{Corollary}
\newcommand{\C}{\mathbb C}
\newcommand{\D}{\mathbb D}
\newcommand{\T}{\mathbb T}
\newcommand{\cl}[1]{\overline{#1}}

\title{A classical strongly regular Morris Swiss cheese\\
\large An affirmative answer to the existential half of Question 1 in
arXiv:1412.6406}
\author{Autonomous mathematical discovery run \texttt{fa\_banach\_001}\\
\texttt{agent\_lane\_02} \; (GPT5.6)}
\date{11 August 2026}

\begin{document}
\maketitle

\begin{center}
\fcolorbox{black}{gray!8}{\parbox{0.93\linewidth}{
\textbf{Status.} Candidate full affirmative answer to the first (existential)
sentence of Question 1, likely valid, pending expert review.  For every
$C>0$ the cheese in Theorem 8.3 can be chosen so that its rational uniform
algebra is in fact strongly regular.  The universal ``must'' sentence and the
two broader questions in the source remain open.}}
\end{center}

\section{Source question and exact scope}

Feinstein, Morley, and Yang prove that for each $C>0$ there is a classical
abstract Swiss cheese $B=((b_n,s_n))$ with outer disk $\cl\D$, with
$\T\subset X_B$, with no nonzero bounded point derivations on $R(X_B)$, and
with
\begin{equation}\label{eq:weighted}
       \sum_{n\geq1}\frac{s_n}{d_n^2}<C,
       \qquad d_n=\operatorname{dist}(B(b_n,s_n),\T).
\end{equation}
The weighted estimate makes
$(f,g)\mapsto\int_{\T}f'(z)g(z)\,dz$ a bounded bilinear form and hence gives a
nonzero derivation into the dual \cite{FMY}.  On source PDF page 18 they ask:

\begin{figure}[ht]
  \centering
  \includegraphics[width=\textwidth]{figures/open_problem_crop.png}
  \caption{Question 1 (and the adjacent Question 2) on page 18 of the locally
  compiled arXiv source release of arXiv:1412.6406.}
\end{figure}

\begin{quote}
\textbf{Question 1.} Let $B$ be the classical abstract Swiss cheese
constructed in Theorem 8.3. Can $R(X_B)$ be regular? Must $R(X_B)$ be
regular?
\end{quote}

This packet proves that the answer to \emph{Can it be regular?} is yes, in the
stronger form ``can be strongly regular.''  It does not claim that every
choice in Theorem 8.3 is regular.

\section{Definitions}

For a compact plane set $K$, let $R(K)$ be the uniform closure on $K$ of the
rational functions with poles off $K$.  At $x\in K$ put
\[
 M_x=\{f\in R(K):f(x)=0\},\qquad
 J_x=\{f\in R(K):f\text{ vanishes on a relative neighborhood of }x\}.
\]
The algebra is \emph{strongly regular} if $\cl{J_x}=M_x$ for every $x$.
For an open disk $D_j=B(c_j,r_j)$ whose closure is contained in $\D$, write
\[
          d_j=\operatorname{dist}(D_j,\T)=1-|c_j|-r_j,
          \qquad
          W(\{D_j\})=\sum_j\frac{r_j}{d_j^2}.
\]

\section{Proof intuition}

Izzo's 2022 construction gives two complementary local ideal estimates.  A
quantitative K\"orner construction forces $\cl{J_x}\supset M_x^3$, while a
localized Wermer construction forces $\cl{M_x^2}=M_x$.  Together they imply
strong regularity.  Izzo arranged the first estimate with a weighted boundary
sum on a square \cite{Izzo}; the first new point here is that a Cayley chart
flattens the unit circle \emph{exactly}.  Because a M\"obius map takes disks to
disks, its half-plane version transfers the boundary construction to the unit
disk with the same type of weighted control.

The resulting cheese need not be classical.  The second new point is a short
fixed-boundary classicalisation lemma.  During the standard merging process,
every aggregate disk has radius at most the sum of the original radii assigned
to it.  A small value of $\sum r_j/d_j^2$ then forces that aggregate to remain
at least half as far from $\T$ as its closest assigned disk.  Thus the outer
unit disk never has to shrink, and the weighted sum increases by at most a
factor nine.  Strong regularity survives because the two ideal estimates used
to obtain it pass to compact subsets.

\section{A weighted unit-disk form of Izzo's construction}

We isolate the disk version needed below.  Its proof records the changes to
Izzo's square argument, rather than treating the change of boundary as a
formal affine transformation (which it is not).

\begin{proposition}\label{prop:disk-izzo}
For every $\varepsilon>0$ there is a sequence of open disks $D_j$ with
$\cl{D_j}\subset\D$ such that, for
\[
                 K=\cl\D\setminus\bigcup_jD_j,
\]
one has $W(\{D_j\})<\varepsilon$ and, for every $x\in K$,
\begin{equation}\label{eq:two-ideals}
 \cl{J_x(R(K))}\supset M_x(R(K))^3,
 \qquad
 \cl{M_x(R(K))^2}=M_x(R(K)).
\end{equation}
Consequently $R(K)$ is strongly regular.
\end{proposition}

\begin{proof}
We use the notation and quantitative K\"orner lemmas of Izzo
\cite[Lemmas 3.1, 3.5, 4.1, and 6.7]{Izzo}.  Those lemmas have two features
that matter here.  Given a center $a$, an error $\eta>0$, and a sufficiently
small scale $\rho$, they produce pairwise disjoint pole disks in
\[
             \rho<|w-a|<2\rho
\]
and a function which vanishes on $|w-a|\leq\rho$ and approximates the relevant
coordinate power.  In the boundary-weighted version the power is three and
\begin{equation}\label{eq:local-weight}
 \sum_Q\frac{r(Q)}{s_a(Q)^2}<\frac{4\delta}{\rho},
\end{equation}
where $s_a(Q)$ is the distance from $Q$ to the two coordinate lines through
$a$.  The small parameter $\delta$ is free.

We first explain the only genuinely new local step, at a point $a\in\T$.
After rotation use the Cayley map
\[
              \Phi_a(z)=\frac{a-z}{a+z}.
\]
It maps $a$ to $0$, $-a$ to infinity, $\D$ to the right half-plane, and
$\T$ to the imaginary axis.  Apply the boundary-weighted local construction
at $0$.  No pole disk meets either coordinate axis (otherwise the left side
of \eqref{eq:local-weight} would be infinite).  Retain the disks in the right
half-plane and discard those in the left half-plane.  Their inverse images
under $\Phi_a$ are disks with closures in $\D$; the discarded poles map
outside $\cl\D$.

On $|w|\leq2\rho$, the inverse Cayley map and its derivative have uniform
upper and lower bounds.  It follows directly from the formula for the image
of a circle under a M\"obius map that, with $Q^*=\Phi_a^{-1}(Q)$,
\begin{equation}\label{eq:mobius-weight}
 \frac{r(Q^*)}{\operatorname{dist}(Q^*,\T)^2}
 \leq A\frac{r(Q)}{\operatorname{dist}(Q,i\mathbb R)^2}
 \leq A\frac{r(Q)}{s_0(Q)^2},
\end{equation}
where the absolute constant $A$ is independent of $a$ and of small $\rho$.
If $h$ is Izzo's local function, set
\[
              G(z)=(z-a)^3h(\Phi_a(z)).
\]
Since $z-a=-2aw/(1+w)$ and $\operatorname{Re}w\geq0$ on $\cl\D$,
\[
 |(z-a)^3-G(z)|
 \leq 8|w^3(1-h(w))|\leq 8\eta.
\]
The apparent singularity at $-a$ is removable: the strengthened K\"orner
function has $h(w)=1+O(w^{-3})$ at infinity.  Its rational approximants after
composition likewise have all poles either in the retained deleted disks or
outside $\cl\D$.  Hence $G$ belongs to the appropriate rational algebra,
vanishes on a relative neighborhood of $a$, and approximates $(z-a)^3$.

Cover $\T$ at the $n$th scale by at most $A_1/\rho_n$ such Cayley
neighborhoods.  Take $\eta_n\downarrow0$, $\rho_n\downarrow0$, and then choose
the free parameters $\delta_n$ so fast that
\[
              \sum_n A A_1\frac{\delta_n}{\rho_n^2}
              <\frac{\varepsilon}{4}.
\]
Equations \eqref{eq:local-weight}--\eqref{eq:mobius-weight} give the desired
weighted budget at the boundary, and the preceding approximation gives
$\cl{J_x}\supset M_x^3$ for every $x\in\T$.  More explicitly, the inverse
Cayley image of $|w|<\rho_n$ contains a Euclidean neighborhood of its center
of radius comparable to $\rho_n$; applying Izzo's local criterion
\cite[Lemma 4.1]{Izzo} converts the covering approximations of
$(z-a)^3$ into the asserted ideal inclusion at every covered point.

For interior points, exhaust $\D$ by the disks
$L_n=\{|z|\leq1-1/n\}$.  Cover $L_n$ by at most
$A_2/\sigma_n^2$ disks of radius $\sigma_n$, with centers in $L_n$ and with
$4\sigma_n<1/n$, and use Izzo's ordinary local lemma at every covering center.
(A maximal separated net gives an absolute $A_2$.)  Each pole disk then has distance at least $1/(2n)$ from
$\T$.  The ordinary radius bound therefore makes the total weighted cost at
stage $n$ at most a fixed constant times $n^2\delta_n/\sigma_n$.  Choosing
$\delta_n$ sufficiently fast gives total cost below $\varepsilon/4$ and gives
the stronger inclusion $\cl{J_x}\supset M_x^2$ at every interior point.
Together with the boundary construction this proves the first inclusion in
\eqref{eq:two-ideals} everywhere.

For the second inclusion, take a countable cover $U_j$ of $\D$ by disks with
$\cl{U_j}\subset\D$.  Apply the scaled Wermer construction in each $U_j$,
choosing its ordinary radius sum below
\[
  2^{-j-2}\varepsilon\,\operatorname{dist}(U_j,\T)^2.
\]
The total added weighted cost is below $\varepsilon/4$.  Hallstr\"om's local
criterion, exactly as in Izzo's proof of his weighted Wermer theorem
\cite[Theorem 6.6]{Izzo}, gives $\cl{M_x^2}=M_x$ at every interior point.
Every boundary point is a peak point, using
$(1+\overline a z)/2$, so the same equality holds on $\T$.

Finally take the union of the two pole families.  Izzo's compact-subset
lemmas \cite[Lemmas 5.1 and 5.2]{Izzo} show that both ideal properties pass to
this intersection.  The second one implies that $M_x^3$ is dense in $M_x$;
the first then yields $\cl{J_x}=M_x$ for every $x$.  Rescaling the three budget
allocations slightly, if necessary, makes their total less than the stated
$\varepsilon$.
\end{proof}

\section{Weighted classicalisation without moving the outer circle}

The following quantitative observation is the step that prevents ordinary
classicalisation from losing $\T$.

\begin{lemma}\label{lem:classical}
Let $A$ be a unit-disk Swiss cheese whose deleted disks
$D_j=B(c_j,r_j)$ satisfy $\cl{D_j}\subset\D$ and
\[
                \theta=\sum_j\frac{r_j}{d_j^2}<\frac16.
\]
Then there is a classical unit-disk Swiss cheese $B$ such that
\[
       \T\subset X_B\subset X_A,
       \qquad
       \sum_{E\in B}\frac{r(E)}{\operatorname{dist}(E,\T)^2}
       \leq 9\theta.
\]
The outer disk remains exactly $\cl\D$.
\end{lemma}

\begin{proof}
Run the standard Feinstein--Heath/Mason classicalisation by combining two
inner disks whenever their closures meet \cite{FMY}.  Keep the usual
allocation bookkeeping: every current inner disk $E=B(c,R)$ is assigned a
disjoint set $I(E)$ of original disks, contains each assigned disk, and
satisfies
\begin{equation}\label{eq:allocation}
                 R\leq\sum_{i\in I(E)}r_i.
\end{equation}
These invariants are preserved because the minimal disk containing two
intersecting disks has radius at most the sum of their radii.  They also pass
to the limit stages in the standard construction.

Put $e=\operatorname{dist}(E,\T)=1-|c|-R$ and
$d_* =\inf_{i\in I(E)}d_i$.  Containment of $D_i$ in $E$ gives
\[
       e\leq d_i\leq e+2R\leq d_*+2R
       \qquad(i\in I(E)).
\]
Consequently \eqref{eq:allocation} gives
\begin{equation}\label{eq:selfconsistent}
       R\leq\theta(d_*+2R)^2,
       \qquad R\leq\sum_i r_i\leq\theta,
\end{equation}
because $0<d_i<1$.  If $R\geq d_*/4$, then
\[
       R\leq36\theta R^2,
       \quad\hbox{so}\quad
       1\leq36\theta R\leq36\theta^2<1,
\]
a contradiction.  Thus $R<d_*/4$ and
\begin{equation}\label{eq:positive-gap}
               e\geq d_*-2R>d_*/2>0.
\end{equation}
Every aggregate disk therefore remains strictly inside $\D$.  The operation
that shrinks or moves the outer disk is never invoked, so the classical output
has outer disk exactly $\cl\D$ and contains $\T$.

For every final disk $E$, equations \eqref{eq:positive-gap} and
$d_i<d_*+2R<3d_*/2$ give
\[
 \frac{R}{e^2}
 <\frac{4R}{d_*^2}
 \leq 9\sum_{i\in I(E)}\frac{r_i}{d_i^2}.
\]
The final allocations are disjoint.  Summing proves the factor-nine bound.
The usual classicalisation inclusion gives $X_B\subset X_A$.
\end{proof}

\section{Main theorem}

\begin{theorem}\label{thm:main}
For every $C>0$ there is a classical abstract Swiss cheese
$B=((b_n,s_n))$ such that
\begin{enumerate}
\item $b_0=0$, $s_0=1$, and $\T\subset X_B$;
\item $R(X_B)$ is strongly regular (hence regular and has no nonzero bounded
point derivations);
\item with $d_n=\operatorname{dist}(B(b_n,s_n),\T)$,
\[
                 \sum_{n\geq1}\frac{s_n}{d_n^2}<C;
\]
\item $R(X_B)$ is not weakly amenable.
\end{enumerate}
Thus the existential sentence of Question 1 in arXiv:1412.6406 has an
affirmative answer.
\end{theorem}

\begin{proof}
Choose
\[
             0<\varepsilon<\min\{1/6,C/9\}
\]
and take the unit-disk Swiss cheese $A$ from
Proposition~\ref{prop:disk-izzo}.  Apply Lemma~\ref{lem:classical} to obtain a
classical $B$ with $X_B\subset X_A$, fixed outer disk, and weighted sum below
$9\varepsilon<C$.

The two ideal conditions in \eqref{eq:two-ideals} pass from $X_A$ to the
compact subset $X_B$ by the same restriction argument used in
\cite[Lemmas 5.1 and 5.2]{Izzo}.  They again imply
$\cl{J_x(R(X_B))}=M_x(R(X_B))$ for all $x\in X_B$.  Thus $R(X_B)$ is strongly
regular.  In particular it is regular and has no nonzero bounded point
derivations.

The Cauchy estimate of Feinstein--Morley--Yang
\cite[Lemma 8.2]{FMY}, together with the weighted sum, gives
\[
 \left|\int_{\T}f'(z)g(z)\,dz\right|
 \leq4\pi C\|f\|_{X_B}\|g\|_{X_B}
 \qquad(f,g\in R_0(X_B)).
\]
There is at least one deleted disk because $X_B\subset X_A\ne\cl\D$.  Choose
$q$ in one of them.  For $f(z)=1/(z-q)$ and $g(z)=z-q$ the integral is
$-2\pi i$.  Hence the associated bounded derivation
$R(X_B)\to R(X_B)^*$ is nonzero, so $R(X_B)$ is not weakly amenable.
All conclusions of Theorem 8.3, with the strengthened regularity conclusion,
now hold.
\end{proof}

\section{Verification, limitations, and novelty status}

The proof has four independent checkpoints.
\begin{enumerate}
\item The Cayley map sends the protected circle exactly to an axis and the
unit disk to one half-plane; the removability at the antipodal point uses the
third-order normalization already present in Izzo's weighted construction.
\item The interior and boundary covering counts are respectively
$O(\sigma^{-2})$ and $O(\rho^{-1})$.  The free K\"orner parameters can be
chosen after the scales, so both weighted series converge with an arbitrarily
small sum.
\item In Lemma~\ref{lem:classical}, the contradiction is strict because
$\theta<1/6$, and the factor nine follows from the two sharp comparisons
$e>d_*/2$ and $d_i<3d_*/2$.
\item Strong regularity is transferred through the two separate inherited
ideal statements, not through an unsupported assertion that strong
regularity is hereditary for arbitrary uniform algebras.
\end{enumerate}

No computation is used in the proof.  The source PDF and the decisive Izzo
paper are included in this packet.  A bounded novelty search covered the
exact wording of Question 1, ``classical Swiss cheese strongly regular not
weakly amenable,'' the source authors, arXiv:1501.04008, arXiv:1503.03785,
arXiv:1507.01779, and arXiv:2211.14684, using the local corpus and exact web
searches through 11 August 2026.  The 2016 paper repeats the broader Question
2; Izzo proves the square weighted theorem but does not state the classical
unit-disk conclusion.  No prior statement of Theorem~\ref{thm:main} was
found.  This is a bounded search, not a claim of exhaustive novelty.

The result answers only the existential ``Can'' sentence of source Question
1.  It does not show that every cheese produced by Theorem 8.3 is regular,
does not answer whether absence of bounded point derivations always implies
regularity, and does not answer whether weak amenability forces triviality.

\paragraph{Human-review recommendation.}
Check first the M\"obius transfer in Proposition~\ref{prop:disk-izzo},
especially removability at $-a$ and the uniform distance comparison
\eqref{eq:mobius-weight}.  Then check that the standard classicalisation
allocation process has exactly the containment and radius invariants used in
Lemma~\ref{lem:classical}.  If those checks pass, accept the existential half
of Question 1 as fully solved.

\begin{thebibliography}{9}
\bibitem{FMY}
J. F. Feinstein, S. Morley, and H. Yang,
\emph{Swiss Cheeses and Their Applications}, arXiv:1412.6406,
Contemporary Mathematics 645 (2015), DOI: 10.1090/conm/645/12928.

\bibitem{Izzo}
A. J. Izzo,
\emph{A sharper Swiss cheese}, arXiv:2211.14684 (revised 2025).

\bibitem{FHMC}
J. F. Feinstein, S. Morley, and H. Yang,
\emph{Abstract Swiss Cheese Space and Classicalisation of Swiss Cheeses},
J. Math. Anal. Appl. 438 (2016), 119--141; arXiv:1503.03785.
\end{thebibliography}

\end{document}
